\section{Manual de operación y referencia de herramientas}

\subsection{Reward Learning Action Timeline}

\begin{table}[H]
\centering
\begin{tabular}{p{0.25\textwidth} p{0.45\textwidth} p{0.2\textwidth}}
\toprule
Etapa & Actividad principal & Responsable \\
\midrule
Preparación de datos & Recopilar ejemplos positive/negative, comparaciones de preferencia & Data engineer + labelers \\
Reward training & Entrenar el classifier/RLM, monitorear la curva de loss & ML engineer \\
Evaluation & Métricas automáticas, human eval, safety probe & Ops + safety reviewer \\
Deployment & KL penalty tuning, logging + rollback & DevOps \\
\bottomrule
\end{tabular}
\caption{Línea de tiempo operativa de reward learning}
\end{table}

\subsection{Herramientas y comandos}

\begin{itemize}
    \item Usar `tensorboard` para monitorear los logits + loss del reward model
    \item Escribir `reward_score` en el experiment tracker (Weights \& Biases / Neptune)
    \item Para el pipeline de RLHF se recomienda usar una combinación de `trl` + `DPO` + `PPO`
\end{itemize}

\begin{knowledgebox}{Ubicación de recursos}
\textbf{Los materiales de reward learning} se concentran en: OpenAI blog, Anthropic interpretability notes, Stanford RL lab releases. Mantenerse suscrito a estas actualizaciones permite ajustar a tiempo la payoff curve.
\end{knowledgebox}

\begin{figure}[H]
\centering
\includegraphics[width=0.92\textwidth,page=7]{08_cs224r_reward_learning_2025.pdf}
\caption{Ejemplo de reward trace dashboard (lecture slide, page 7)}
\end{figure}

\subsection{Resumen de la sección}

El manual de operación enfatiza que cada etapa necesita un responsable que avance en paralelo, sobre todo porque reward training y evaluation no pueden desacoplarse.

